\documentclass[11pt,a4paper]{article}

\usepackage[utf8]{inputenc}
\usepackage[T1]{fontenc}
\usepackage{lmodern}
\usepackage[margin=1in]{geometry}
\usepackage{amsmath,amssymb,amsthm,mathtools}
\usepackage{graphicx}
\usepackage{booktabs}
\usepackage{hyperref}
\usepackage{xcolor}
\usepackage{enumitem}

\hypersetup{
  colorlinks=true,
  linkcolor=blue!50!black,
  citecolor=blue!50!black,
  urlcolor=blue!50!black
}

\title{\bfseries Probing, scattering, and dimensional extensions\\
to the Systroph\=e framework\\
\large Five further modules (Phases 21--25)}

\author{Christian Knopp\\\texttt{cknopp@gmail.com}}
\date{2026-05-11}

\begin{document}
\maketitle

\begin{abstract}
We add five further modules to the Systroph\=e framework covering
probe diagnostics, wave scattering, numerical-relativity hand-off,
dark-matter coupling, and higher-dimensional embedding:
(1) tidal-force scalar across CTC bands, including geodesic-deviation
evolution and a safety-radius API;
(2) Klein-Gordon scattering with transmission/reflection coefficients
across chronology horizons, including a superradiance probe;
(3) NR initial-data export (ADM + BSSN) with constraint-residual
diagnostics and 3D Cartesian extrusion;
(4) dark-matter scalar coupling treating the cylinder as an ultra-light
DM soliton anchor;
(5) 5D Kaluza-Klein embedding showing that motion in a compact extra
dimension can convert a 4D CTC into a 5D spacelike path.
Seventy-one new tests across the five modules, all passing.
Code at \texttt{github.com/Zynerji/systrophe} (v0.19.0+).
\end{abstract}

\section{Tidal forces across CTC bands (Phase 21)}

A probe near a chronology horizon experiences strong tidal forces.
The dominant radial tidal scalar is
\[
T(r) = -\tfrac{1}{2}\,F''(r),
\]
which oscillates log-periodically and diverges as $r$ approaches
each F = 0 surface.

The module \texttt{tidal\_forces.py} provides:
\begin{itemize}
\item \texttt{riemann\_scalar\_radial}: $T(r)$;
\item \texttt{tidal\_stretch\_at\_radius}, \texttt{tidal\_squeeze\_at\_radius};
\item \texttt{geodesic\_deviation\_evolution}: Jacobi equation integration;
\item \texttt{multi\_band\_tidal\_signature}: $T$ at the midpoint of each CTC band;
\item \texttt{safety\_radius\_for\_probe}: where $|T|$ exceeds a material strength.
\end{itemize}

For omega = R = 1, $|T|$ stays small in the first CTC band but grows
log-periodically toward the second, then larger horizons -- a probe
made of any finite material has a definite stay-out radius.

\section{Klein-Gordon scattering (Phase 22)}

A massive scalar $\phi$ on the LP exterior obeys
$(\Box - m^2)\phi = 0$. Decomposing $\phi = e^{-i\omega t}e^{in\phi}\chi(r)$,
the radial equation reduces to a Schrödinger-like form with effective
potential
\[
V_{\rm eff}(r;\omega,n,m) = \frac{\omega^2}{F_{\rm eff}(r)} - m^2 - \frac{n^2}{|L|}.
\]

The module \texttt{kg\_scattering.py} provides:
\begin{itemize}
\item \texttt{effective\_potential};
\item \texttt{transmission\_coefficient}: WKB $|T|^2 = \exp\!\left(-2\int\sqrt{|V_{\rm eff}|}\,dr\right)$;
\item \texttt{reflection\_coefficient}: $|R|^2 = 1 - |T|^2$;
\item \texttt{transmission\_at\_multiple\_CH}: cumulative transmission;
\item \texttt{superradiance\_test}: Kerr-analog amplification condition $\omega < n\,\Omega_H$;
\item \texttt{bound\_states\_between\_CHs}: discrete spectrum in CTC bands.
\end{itemize}

Each chronology horizon acts as a partial barrier; cumulative
transmission decreases geometrically across multiple CHs. Modes with
$\omega < n\,\Omega_H$ are amplified -- the LP analog of Kerr
superradiance.

\section{NR initial-data export (Phase 23)}

The module \texttt{nr\_initial\_data.py} provides constraint-checked
hand-off to numerical-relativity codes (Einstein Toolkit, GRChombo,
BAM):
\begin{itemize}
\item \texttt{adm\_profile\_1d}: lapse, shift, spatial metric, extrinsic curvature on a 1D radial profile;
\item \texttt{hamiltonian\_constraint\_residual} and
  \texttt{momentum\_constraint\_residual}: vacuum-constraint residuals;
\item \texttt{bssn\_decomposition}: conformal factor $\chi$, conformal
  metric $\tilde\gamma$, traceless $\tilde A$;
\item \texttt{extrude\_to\_3d\_cartesian}: 3D grid for NR import;
\item \texttt{export\_initial\_data\_json}: JSON serialization;
\item \texttt{cauchy\_horizon\_warning\_for\_foliation}: detect lapse breakdown.
\end{itemize}

This makes the Systroph\=e framework directly importable into
production NR codes. The Cauchy-horizon warning ensures that the
foliation choice is safe before evolution starts.

\section{Dark-matter scalar coupling (Phase 24)}

For an ultra-light dark-matter scalar ($m_{\rm DM} \sim 10^{-22}$ eV
fuzzy-DM regime), the LP cylinder anchors a soliton-like coherent
field. We model:
\[
\rho_{\rm DM}(r) = \frac{\rho_c}{\left(1 + (r/r_{\rm sol})^2\right)^8}
\]
with $r_{\rm sol}$ set by the Compton wavelength of $m_{\rm DM}$ and
the cylinder mass.

The module \texttt{dm\_scalar\_coupling.py} provides:
\begin{itemize}
\item \texttt{compton\_wavelength};
\item \texttt{DM\_density\_profile\_around\_cylinder};
\item \texttt{DM\_drag\_on\_orbits}: orbital-frequency shift;
\item \texttt{DM\_superradiance\_growth\_rate}: Kerr-analog;
\item \texttt{effective\_alpha\_with\_DM\_pressure};
\item \texttt{DM\_induced\_CTC\_lifetime}: DM-noise decoherence time.
\end{itemize}

For realistic ULDM densities, the DM-noise CTC lifetime exceeds
the Hubble time -- so even pessimistic DM scenarios do not erase
existing Systroph\=e CTCs.

\section{5D Kaluza-Klein embedding (Phase 25)}

Embed the 4D LP metric in 5D with a compact $\xi$ of radius $L_\xi$:
\[
ds_5^2 = ds_{\rm LP}^2 + L_\xi^2\,d\xi^2,\qquad \xi\sim\xi+2\pi.
\]
KK reduction yields a $U(1)$ gauge field plus a dilaton from
$g_{\xi\xi}$. A worldline that moves rapidly in $\xi$ acquires a
positive 5D contribution that can offset a negative 4D one --
\emph{thereby converting a 4D CTC into a 5D spacelike path}.

The module \texttt{kk\_embedding.py} provides:
\begin{itemize}
\item \texttt{five\_d\_metric}: $5\times 5$ tensor;
\item \texttt{kk\_tower\_masses}: $m_n = n/L_\xi$;
\item \texttt{kk\_reduced\_alpha}: $\alpha\sqrt{1 + (R/L_\xi)^2}$;
\item \texttt{xi\_circulation\_evades\_CTC}: minimum $d\xi/d\phi$ to
  achieve a spacelike 5D path inside a 4D CTC band;
\item \texttt{kk\_monopole\_charge}: $U(1)$ charge from non-trivial fiber;
\item \texttt{compactification\_radius\_from\_observations}: bound $L_\xi$;
\item \texttt{five\_d\_geodesic\_drift\_in\_xi}: per-revolution $\xi$ drift.
\end{itemize}

\textbf{Headline}: in any 4D CTC band ($F < 0$), the minimum
$d\xi/d\phi$ for a 5D-spacelike escape is
$\sqrt{-F}/L_\xi$. Sub-Compton $L_\xi$ requires impractical $\xi$
velocities; macroscopic $L_\xi$ would already have been observed
in collider experiments. So KK escape is theoretically possible but
phenomenologically gated.

\section{Implications}

The five modules extend the framework to:
\begin{enumerate}
\item \textbf{Probe physics}: tidal forces give a probe-survivability map.
\item \textbf{Wave scattering}: each chronology horizon is a quantum
  potential barrier with explicit transmission.
\item \textbf{NR pipeline}: initial-data is exportable to Einstein
  Toolkit / GRChombo with constraint residuals tagged.
\item \textbf{Cosmology / DM}: ULDM coupling is weak enough that
  existing CTCs survive Hubble-time DM noise.
\item \textbf{Higher dimensions}: KK extra-dim motion is a
  potential CTC escape route, gated by observational bounds on
  $L_\xi$.
\end{enumerate}

\section{Reproduction}

\begin{center}
\begin{tabular}{ll}
\toprule
Phase & Source file \\
\midrule
21. Tidal forces & \texttt{tidal\_forces.py} \\
22. KG scattering & \texttt{kg\_scattering.py} \\
23. NR initial data & \texttt{nr\_initial\_data.py} \\
24. DM scalar coupling & \texttt{dm\_scalar\_coupling.py} \\
25. 5D KK embedding & \texttt{kk\_embedding.py} \\
\bottomrule
\end{tabular}
\end{center}

All 71 new tests pass; combined Systroph\=e suite at 813 passing tests.

\bibliographystyle{plain}
\begin{thebibliography}{9}

\bibitem{Knopp2026Systrophe}
C.~Knopp, Systrophe whitepaper I (classical GR), 2026.

\bibitem{Knopp2026Ext1}
C.~Knopp, Systrophe extensions paper (Phases 1--10), 2026.

\bibitem{Knopp2026Ext2}
C.~Knopp, Systrophe extensions paper (Phases 11--15), 2026.

\bibitem{Knopp2026Ext3}
C.~Knopp, Systrophe extensions paper (Phases 16--20), 2026.

\bibitem{Baumgarte2010}
T.~Baumgarte and S.~Shapiro, Numerical Relativity, Cambridge, 2010.

\bibitem{Hu2000}
W.~Hu, R.~Barkana, A.~Gruzinov, Phys. Rev. Lett. 85 (2000) 1158.

\bibitem{KK1921}
T.~Kaluza, Sitz. Preuss. Akad. Wiss. (1921) 966; O.~Klein, Z. Phys.
37 (1926) 895.

\end{thebibliography}

\end{document}
